\section{What the SVD Cannot Tell You}
\label{sec:svd-limits}

The SVD is the best document about this chip and it is still not enough. Three classes of fact are
outside it, and every one of them cost this project real time.

\subsection{Class 1: The Board}

An SVD describes a \emph{part}, and a board is a part plus wires:

\begin{itemize}
    \item which pin the LED is on, and whether it is active low;
    \item which two pins carry a bus, and in which order;
    \item what external pull-ups exist;
    \item which pin gates power to a sensor.
\end{itemize}

None of that is in XML, which is why chapter~4 ends with a table of contradictions between the
vendor's board documents and the board. The general rule: \textbf{chip facts from the SVD, board
facts from a measurement, and the measurement written next to the code that depends on it.}

\subsection{Class 2: The Aliases}

This part's peripheral map has two windows: a secure alias at \codeinline{0x4D00\_0000} and a
non-secure one at \codeinline{0x5D00\_0000}, for every peripheral that is not fixed to one domain.
The SVD lists both, and it cannot list which one \emph{your image} may use — that depends on the
security attribution the bootloader and the linker applied.

Measured, on the MVP, from a non-secure image:

\begin{minted}{text}
0x4D00_0000 (secure alias, from non-secure code)  -> reads as zero
0x5D00_0000 (non-secure alias) FAULTSTATUS (0x24) -> bus error
0x5D00_0000 (non-secure alias) the rest           -> works
\end{minted}

Three different behaviours from one block, none of them described by the register file. The practical
consequence is a rule this project now follows for any peripheral with two aliases: \textbf{probe
reachability with a read-back of a harmless register before configuring anything}, and report which
alias answered.

\subsection{Class 3: Timing and Sequencing}

An SVD says which bits exist. It does not say:

\begin{itemize}
    \item that \codeinline{EN} must be written and \codeinline{DISABLING} polled before configuring;
    \item that a converter's \codeinline{SYNCBUSY} must clear between writes;
    \item that a timer's compare register is buffered and takes effect on the next period;
    \item that an oscillator must be forced on and waited for.
\end{itemize}

Those are in the reference manual's prose, and the only reliable way to learn them is to watch a
peripheral misbehave and then find the paragraph. What helps is knowing they exist: a peripheral that
ignores a correct-looking configuration is almost always waiting for a sequence, not a bit.

\subsection{The Order to Read Things In}

\begin{enumerate}
    \item \textbf{Generated register accessors} — for what exists, where, and with which fields.
    \item \textbf{Field descriptions} — for what a value means; they are the SVD's own text.
    \item \textbf{Reference manual prose} — for sequencing and timing, searched by register name.
    \item \textbf{The board, with a probe} — for every wire, polarity and address.
\end{enumerate}

Skipping to step 4 is expensive but sometimes correct; skipping step 1 is never correct, and it is the
most common mistake in a project that started from someone's blog post.
